\section{Our Graph-Based Artificial Intelligence Taxonomy}
\label{section:12-graph-based-ai}
The \gai Benchmark Framework defines Graph-Based Artificial intelligence (\gai) as a taxonomy of intelligence algorithms and machine learning models that learn and reason over information represented in graph data structures.

    
    \begin{figure}[h!]
        \begin{center}
            \includesvg[width=\columnwidth]{SVGs/gai_taxonomy}
            \caption{The \gai Taxonomy with example tasks}
            \label{figure:gai-taxonomy}
        \end{center}
    \end{figure}



Prior efforts to define a taxonomy of \gai focus primarily on \textit{Graph Neural Network}~\cite{Smith2021, Zhou2020a}, \textit{Graph Mining}~\cite{Chen2022a} or specific problems like signal processing~\cite{Xia2021}.
We extend these taxonomies to include:
\begin{itemize}
    \item \textbf{Graph Problem Solving \ac{GPS}.} Classic AI tasks that formulate the world as a graph, including navigation and critical path analysis. 
    \item \textbf{Graph Mining (\ac{GM}).} Analyze the structural characteristics of a graph, including spatial pattern matching and community detection. 
    \item \textbf{Graph Reasoning (\ac{GR}).} AI tasks that use graphs for inference, like common sense reasoning and knowledge-graph question answering. 
    \item \textbf{Graph Abstraction (\ac{GA}).} Transforming graphs into other structures for downstream tasks, for example, generating embeddings with Node2Vec~\cite{Grover2016} or DeepWalk~\cite{Perozzi2014}.
    \item \textbf{Graph Prediction (\ac{GP}).} Learning directly on graphs typically using \ac{GNN}s, including knowledge graph completion and anomaly detection. 
    \item \textbf{Graph Transparency (\ac{GT}).} Uses graphs to improve generative models' transparency, explainability and interpretability. 
\end{itemize}

The taxonomy contextualizes \gai as part of the existing AI landscape and provides a few examples \gai tasks that sit within each \gai category. 
The \gai taxonomy anchors the \gai Benchmark Framework. 
The \gai problem areas guide the development of the benchmark problems within the framework and emphasize the diversity of tasks that the hardware and software systems that run \gai need to support. 

\osullikomment{Add table into Preliminaries}